\section{Observability design}
A pipeline that only works is not enough. We must see when it stops working, and why. Our design follows
three questions: \emph{is data flowing?}, \emph{is it correct?} and \emph{is it on time?} Every stage answers
them with logs, metrics and alerts.

\subsection{Structured logging}
Every Python component (simulators, both Spark jobs, Airflow tasks and the API) writes one JSON object per
line to standard output with the same core keys: \code{ts}, \code{level}, \code{service}, \code{stage}
(ingestion, processing, storage, serving or orchestration) and \code{event}. Extra keys give context, such as
the micro-batch id, the Airflow run id or the simulated date. Spark's own Java logs use Spark's JSON layout,
and the Kafka client logs pass through the same JSON formatter. Because all stages use the same keys, one
\code{jq} filter can follow a problem through the whole pipeline. This is a real line from the speed layer:

\noindent\begin{minipage}{\linewidth}
\begin{Verbatim}[fontsize=\scriptsize,frame=single,framesep=2mm,rulecolor=\color{black!30}]
{"ts": "2026-09-30T08:39:41.660+00:00", "level": "INFO", "service": "spark-speed-layer",
 "stage": "processing", "event": "vehicle_state_batch", "batch_id": 173, "rows": 100,
 "valid": 99, "dead_letters": 1, "vehicles_updated": 20, "idle_alerts_opened": 0,
 "idle_alerts_resolved": 0, "stream_time": "2026-09-02 18:11:36", "e2e_latency_s": 2.05,
 "duration_ms": 1398}
\end{Verbatim}
\end{minipage}

\subsection{Metrics}
Prometheus scrapes every stage every 10 seconds. Table~\ref{tab:metrics} lists the most useful metrics.

\begin{table}[!htb]
\centering\small
\caption{Main metrics, grouped by pipeline stage.}
\label{tab:metrics}
\begin{tabularx}{\textwidth}{@{}P{2.2cm} P{5.4cm} L@{}}
\toprule
\textbf{Stage} & \textbf{Metric} & \textbf{Why we measure it} \\
\midrule
Ingestion & events sent per partition, send errors, delivery latency, time of the last cost file & Is data reaching Kafka? Did the partner file arrive? \\
Processing & input rows per second, micro-batch duration, state rows (from a Spark query listener) & Is Spark keeping up with its trigger? \\
Processing & offsets behind the latest offset (lag), rows dropped by the watermark & Is the speed layer falling behind or losing late data? \\
Storage & end-to-end latency from Kafka to PostgreSQL, rows written, replayed batches & Are the tables fresh, and are writes exactly once? \\
Orchestration & Airflow task results and scheduler heartbeat (StatsD), batch days behind & Is the batch layer running every day? \\
Serving & API request rate and latency, vehicles by status, idle minutes, loss-making vehicles & Is serving healthy, and what is the business state? \\
\bottomrule
\end{tabularx}
\end{table}

\subsection{Tracing and lineage}
Full distributed tracing (for example with Jaeger) needs more memory than our laptop has, so we built
light-weight tracing and lineage from ids that already exist. Every event has an \code{event\_id}, which is
also sent as the Kafka header \code{trace\_id}. The master dataset stores the Kafka topic, partition, offset
and append time of every event, so any number in a report can be traced back to the exact Kafka records.
\code{ops.speed\_batches} records every committed micro-batch, \code{ops.batch\_files} every partner file
with its checksum, and \code{ops.pipeline\_runs} every Airflow step with its inputs, outputs and timing.
Each batch view row also carries the Airflow run id that produced it.

\subsection{Alerts and health checks}
\label{sec:alerts}
Prometheus evaluates 14 rules. Alertmanager groups the alerts and sends them to a webhook in our API, which
writes them as JSON logs and stores them in \code{ops.alert\_notifications}. Airflow task failures are
stored in the same table by a failure callback, so there is one place to see all problems.
Table~\ref{tab:alerts} shows the main rules. Some alerts are hidden when their cause is already known. For
example, if the Spark container is down, the alert ``speed layer stale'' is not sent again.

\begin{table}[!htb]
\centering\small
\caption{Main alert rules (real seconds, simulated time in brackets).}
\label{tab:alerts}
\begin{tabularx}{\textwidth}{@{}P{4.7cm} L P{1.5cm}@{}}
\toprule
\textbf{Alert} & \textbf{Condition} & \textbf{Severity} \\
\midrule
NoTelemetryIngested & no event reached Kafka for 1 minute (2 hours) & critical \\
SpeedLayerStale & the speed layer wrote nothing for more than 60 s & critical \\
StreamingQueryDown, ServiceDown & a Spark query stopped, or a service does not answer & critical \\
KafkaConsumerLagHigh & a query is more than 2000 offsets behind & warning \\
EndToEndLatencyHigh & Kafka to PostgreSQL takes more than 20 s & warning \\
InvalidEventRateHigh & more than 5\% of events are invalid & warning \\
BatchLayerBehind & more than one finished day is not processed & critical \\
ExpenseFeedMissing & no partner file for 1.5 simulated days & critical \\
VehicleIdleTooLong & a vehicle is idle for more than 90 simulated minutes & warning \\
VehiclesBecomingUnprofitable & yesterday's report found vehicles that are becoming unprofitable & info \\
\bottomrule
\end{tabularx}
\end{table}

\textbf{Failure drill.} To test detection, we stopped the stream simulator for 215 real seconds (about
seven simulated hours) and timed each alert.
\DrillText{}
